\chapter[INTRODUÇÃO]{INTRODUÇÃO}

\section{Contextualização}

O Espírito Santo responde pela maior parte da produção e da exportação brasileira de rochas ornamentais, atividade que sustenta uma cadeia industrial extensa, da lavra ao beneficiamento \cite{stone_economy_bahiense2023}. Ao fim dessa cadeia, cada chapa polida passa por uma etapa de inspeção visual em que um operador decide se a peça segue para o cliente, é reclassificada ou é descartada. Essa decisão determina o preço da chapa e, agregada, uma fração relevante da margem da empresa \cite{stone_inspection_vidal2013}.

A inspeção é feita a olho nu, sob iluminação variável e em ritmo de produção. O operador avalia a superfície da chapa procurando trincas, manchas, oxidações e descontinuidades que comprometam o uso comercial da peça. O problema não está na acuidade visual de quem inspeciona, mas em algo anterior: não existe um critério escrito que separe uma feição natural da rocha de um defeito comercial. A fronteira entre as duas categorias depende da litologia, do cliente, do lote e do acabamento contratado e, na prática, depende do inspetor \cite{stone_subjectivity_castro2018}.

Essa ausência tem uma consequência que costuma passar despercebida: o critério existe, mas vive implícito na experiência de cada operador. Ele não pode ser consultado, discutido, comparado entre turnos nem auditado por um cliente que conteste uma reclassificação. Duas chapas equivalentes podem receber destinos diferentes sem que haja registro de qual regra produziu cada decisão.

A visão computacional é uma resposta natural a esse tipo de problema, e o setor já registra tentativas de automação da inspeção de superfícies pétreas \cite{stone_tereso2020, stone_lopez2010}. A dificuldade é que as abordagens supervisionadas convencionais exigem grandes volumes de imagens anotadas manualmente, e anotar anomalias sutis em chapas de rocha é caro, lento e, pelo mesmo motivo descrito acima, subjetivo. Anotar milhares de imagens apenas transfere a arbitrariedade do chão de fábrica para o conjunto de treinamento. O efeito já foi observado sobre o mesmo conjunto de dados utilizado neste trabalho: em um segmentador treinado com anotação manual de defeitos em quatro classes, parte das detecções contabilizadas como falsos positivos eram defeitos reais, deixados fora da referência pelo critério de anotação adotado \cite{fracstoneai_oliveira2026}. O critério de quem anota passa a governar a própria métrica.

\section{Problema de Pesquisa}

O problema que este trabalho ataca é a arbitrariedade não registrada do critério de marcação, mais do que a dificuldade técnica de detectar anomalias. Não existe uma resposta absolutamente correta a ser descoberta: o que é defeito em uma litologia pode ser exatamente o padrão estético valorizado em outra. Admitir isso muda o que se pode legitimamente prometer. Um sistema automatizado não elimina a arbitrariedade da fronteira; o que ele pode fazer é retirá-la da cabeça do inspetor e colocá-la em um parâmetro inspecionável.

Formula-se, assim, a seguinte pergunta de pesquisa: \textit{é possível materializar o critério de marcação de anomalias em chapas de rochas ornamentais como um conjunto de parâmetros explícito, versionável e específico por litologia? E esse critério parametrizado produz marcações mais alinhadas ao julgamento humano do que um critério único aplicado a todas as rochas?}

\section{Hipótese}

A hipótese que orienta o trabalho é:

\begin{description}
  \item[H1] Um critério de marcação parametrizado por litologia produz marcações mais alinhadas ao julgamento humano, medido em auditoria cega, do que um critério único e global.
\end{description}

Julgamento humano, aqui, é a resposta de uma avaliadora que não participou da calibração à pergunta ``você marcaria esta região como defeito?'', feita sobre cada marcação sem que ela saiba de que configuração a marcação veio. Do julgamento saem duas medidas, a precisão e o \textit{recall} relativo, e a regra que decide a hipótese quando as duas apontam em sentidos opostos foi fixada antes de qualquer julgamento (Seção \ref{sec:decisao_h1}).

A hipótese trata do critério que gera as marcações, e não do modelo que depois aprende a reproduzi-las. A pergunta seguinte, se segmentadores especializados por litologia superam um único modelo treinado sobre todas e a partir de quantas imagens a especialização compensa o custo de manter um modelo dedicado, depende de as marcações já estarem calibradas e avaliadas, e fica registrada como continuação do trabalho (Seção \ref{sec:futuros}).

\section{Objetivos}

\subsection{Objetivo Geral}

Desenvolver e avaliar o ARIA (Análise e Reconhecimento Inteligente de Anomalias), um \textit{pipeline} hierárquico de visão computacional que materializa o critério de marcação de anomalias superficiais em chapas de rochas ornamentais como um conjunto explícito e versionável de parâmetros por litologia, e verificar se essa parametrização produz marcações mais alinhadas ao julgamento humano do que um critério único global.

\subsection{Objetivos Específicos}

\begin{enumerate}
  \item Selecionar, caracterizar e pré-processar um conjunto de dados público de imagens de chapas polidas de rochas ornamentais, descrevendo sua composição, o desbalanceamento entre litologias e o particionamento adotado.
  \item Reproduzir, sobre as partições originais desse conjunto, o classificador de litologias de trabalho anterior, que identifica a rocha e seleciona a configuração de marcação correspondente.
  \item Definir um protocolo reprodutível de calibração que determine, para cada litologia, quais sondas textuais entram no conjunto e qual limiar de confiança cada uma recebe, de modo que a regra que produz o limiar seja a mesma para todas as litologias.
  \item Implementar o estágio Professor com um modelo de fundação guiado por texto, capaz de gerar polígonos de anomalia sem anotação humana prévia.
  \item Construir um protocolo de avaliação cega, em que avaliadores humanos julgam as marcações sem saber de que configuração vieram, para servir de referência quantitativa da avaliação.
  \item Avaliar experimentalmente a hipótese H1, comparando o critério parametrizado por litologia ao mesmo critério derivado do conjunto reunido de litologias.
\end{enumerate}

\section{Justificativa}

Do ponto de vista industrial, tornar explícito o critério de marcação tem valor independente da acurácia alcançada. Um arquivo de parâmetros pode ser versionado, comparado entre unidades produtivas, negociado com um cliente e auditado quando uma classificação é contestada, o que um critério guardado na experiência de um operador não permite. A automação da marcação é consequência; a contribuição central é a explicitação.

Do ponto de vista acadêmico, o trabalho se situa no encontro de duas linhas ativas. A primeira é o uso de modelos de fundação como geradores de rótulos para domínios sem \textit{ground truth}, o que dispensa a anotação manual em larga escala. A segunda é o aprendizado a partir de rótulos ruidosos, campo que estuda justamente o que acontece quando a supervisão vem de um modelo e não de um humano. O trabalho contribui com um elemento pouco explorado nessa interseção: o condicionamento do \textit{prompt} por classe do domínio, aqui por litologia, como forma de controlar a qualidade do rótulo gerado, e o registro desse condicionamento como artefato auditável.

Por fim, o uso de um conjunto de dados público \cite{dataset_araujo2023} garante que todo o protocolo possa ser reproduzido por terceiros sem acordo com nenhuma empresa, o que é incomum em trabalhos de inspeção industrial.

\section{Organização do Trabalho}

O restante deste documento está organizado da seguinte forma. O Capítulo 2 apresenta a fundamentação teórica, cobrindo a inspeção de qualidade no setor de rochas ornamentais, redes neurais convolucionais aplicadas à classificação de litologias, modelos de fundação com segmentação guiada por texto, a arquitetura Professor--Aluno com pseudo-rotulagem, a escolha do limiar sob rótulos ruidosos, a segmentação em tempo real, as métricas de avaliação e os trabalhos relacionados contra os quais a contribuição se posiciona. O Capítulo 3 detalha a metodologia: o conjunto de dados, a arquitetura do \textit{pipeline}, o protocolo de calibração, a regra que produz o limiar, a varredura de limiares fora de linha, as ferramentas de seleção e calibração, a avaliação por auditoria cega e o desenho do experimento, com o critério que decide a hipótese. O Capítulo 4 apresenta os resultados obtidos: as verificações que sustentam o método, a reprodução do classificador, a calibração das litologias de maior volume e o experimento. O Capítulo 5 traz as considerações finais, o escopo desta versão e os desdobramentos deixados como trabalho futuro. O Apêndice A lista os limiares de todos os pares de litologia e sonda calibrados, e o Apêndice B mostra, para cada litologia da faixa A, as marcações no limiar da regra e no limiar do autor sobre as mesmas chapas.
